\begin{abstract}
Multi-agent LLM systems are commonly evaluated only by their final answer, even though a mechanism must first induce agents to choose informative and costly reasoning.  We introduce \emph{AEIC-MD}, a mechanism-mediated deliberation framework in which Assumption--Evidence--Inference--Conclusion (AEIC) reports provide the observable interface for a public leader rule and agents choose among low-effort, structured AEIC-effort, verification, and imitation actions.  The central empirical question is therefore not whether an offline score can retrospectively rank responses, but whether a revealed rule changes the agents' \emph{actual} effort choices.

Across three QA domains and 450 paired formal questions, a Static-AEIC rule raises the high-effort choice rate from 0.46\% to 56.93\% and decision quality from 68.00\% to 82.22\% (+14.22pp; Holm-adjusted McNemar $p=1.23\times10^{-11}$).  An outcome-only reward changes neither quality (67.33\%, $p=0.780$) nor the main conclusion, separating structured verifiability from generic reward.  Effort responds monotonically to revealed reward and inversely to effort cost in each domain.  Controlled information interventions confirm that AEIC exposes behavior under evidence, social, copying, and counterevidence conditions.  The quality gain also recurs under two additional seeds (+14.67pp; $p=2.40\times10^{-9}$) and remains stable with three to five agents.  However, it does not reproduce on a held-out GLM backend (69.33\% vs.\ 70.67\%, $p=0.791$), establishing a clear boundary rather than a universal claim.  These results position AEIC-MD as a verifiable-effort mechanism for multi-agent deliberation, while identifying backend-specific contract calibration as the key remaining challenge.

\keywords{Multi-agent LLM systems \and Mechanism design \and Incentive compatibility \and Structured reasoning \and Deliberation}
\end{abstract}
